Here is the complete, compiled LaTeX document for the entire Appendices chapter.

I have restored all the mathematical equations, variables, and algorithmic loops that were lost when you copied the plain text, and I have included a standard preamble so you can paste this directly into a blank Overleaf project and compile it immediately.

```latex
\documentclass[12pt]{article}

% --- REQUIRED PACKAGES ---
\usepackage[utf8]{inputenc}
\usepackage{amsmath}
\usepackage{amssymb}
\usepackage{algorithm}
\usepackage{algorithmic}
\usepackage{booktabs}
\usepackage{float}
\usepackage[margin=1in]{geometry}

\begin{document}

\section*{Appendices}

\section*{Appendix A: Denoising Autoencoder (DAE) Architecture and Objective Formulations}

This section details the formal mathematical constructs and specific network topology for the application of the Denoising Autoencoder model utilized in the banking fraud project.

\subsection*{A.1 Tabular Swap Noise Augmentation}
To preserve the univariable marginal feature distributions while destroying joint multi-feature correlations, online data augmentation is applied via column-wise Tabular Swap Noise. For each feature dimension $j$, a stochastic mask variable $m_j \sim \text{Bernoulli}(\alpha)$ is sampled, where the corruption hyperparameter $\alpha$ is fixed at 0.15 (15\% corruption probability).

The mathematical transformation defining the swap noise mapping is:
\begin{equation}
    \tilde{x}_{j} = (1 - m_{j})x_{j} + m_{j} \cdot x_{j}^{*}
\end{equation}
Where $x_{j}^{*}$ represents an attribute value sampled at random from the empirical marginal distribution of the same $j$-th feature channel across an independent transaction record.

\subsection*{A.2 DAE Objective Function}
The network is optimized using a target-isolated Mean Squared Error (MSE) loss function. While the encoder receives the corrupted vector $\tilde{x}$, the loss is computed exclusively between the decoder's reconstructed output $\hat{x}$ and the clean, pristine input vector $x$:
\begin{equation}
    \mathcal{L}_{DAE}(\phi,\theta) = \frac{1}{53}\sum_{j=1}^{53}(x_{j} - \hat{x}_{j})^{2}
\end{equation}

\subsection*{A.3 Network Topology}
The DAE utilizes a deep, symmetric deterministic configuration built with fully connected dense layers, normalized with Batch Normalization, and activated via LeakyReLU (negative slope parameter 0.20).

\begin{itemize}
    \item \textbf{Encoder Stack:} 53 features $\rightarrow$ 128 units $\rightarrow$ 64 units $\rightarrow$ Latent Bottleneck $z$
    \item \textbf{Latent Bottleneck ($z$):} 12 coordinates
    \item \textbf{Decoder Stack:} Latent Bottleneck $z$ $\rightarrow$ 64 units $\rightarrow$ 128 units $\rightarrow$ 53 units (Linear Activation)
\end{itemize}

\vspace{2em}

\section*{Appendix B: Spiking Neural Network (SNN) Mathematical Derivations}

\subsection*{B.1 Leaky Integrate-and-Fire (LIF) Dynamics}
The continuous-time differential equation defining the subthreshold membrane potential $V(t)$ for the LIF neuron is formulated as:
\begin{equation}
    \tau\frac{dV(t)}{dt} = -(V(t) - V_{rest}) + RI(t)
\end{equation}
For implementation within a discrete-time machine learning framework, this is updated iteratively via the explicit Euler method. The subthreshold membrane potential $U[t]$ is calculated as:
\begin{equation}
    U[t] = \beta U[t-1] + I_{syn}[t]
\end{equation}
Where $\beta$ denotes the decay factor and $I_{syn}[t]$ is the weighted synaptic input from presynaptic spikes.

\subsection*{B.2 Surrogate Gradient Descent (ArcTan)}
To permit gradient flow through the non-differentiable Heaviside step function during Backpropagation Through Time (BPTT), an ArcTan (ATan) surrogate gradient is employed. The derivative for the backward pass is strictly computed as:
\begin{equation}
    \frac{\partial S}{\partial U} = \frac{1}{\pi} \frac{1}{1 + \left( \frac{\pi U \alpha}{2} \right)^2}
\end{equation}
Where the hyperparameter $\alpha$ controls the steepness of the gradient approximation.

\newpage

\section*{Appendix C: Graph Neural Network (GNN) Algorithmic Procedures}

\begin{algorithm}[H]
\caption{Quantum Graph Neural Network (QGNN) Forward Pass and Training}
\begin{algorithmic}[1]
\REQUIRE Graph $\mathcal{G}=(\mathcal{V},\mathcal{E})$, features $\{x_u\}$, labels $\{y_u\}$, VQC layers $L$, qubits $n_q$
\ENSURE Trained model parameters $\Theta^*$
\STATE Initialize $W_{enc}, \{W^{(l)}, b^{(l)}\}_{l=1}^L$, MLP parameters
\FOR{epoch $= 1, 2, \dots, E$}
    \FOR{each node $u \in \mathcal{V}$ in mini-batch}
        \STATE $\theta_u \leftarrow [\cos(\pi W_{enc}x_u), \sin(\pi W_{enc}x_u)]$
        \STATE $v_u^{(0)} \leftarrow \theta_u$
        \FOR{$l = 1$ to $L$}
            \STATE $v_u^{(l)} \leftarrow \tanh(W^{(l)}v_u^{(l-1)} + b^{(l)})$
        \ENDFOR
    \ENDFOR
    \FOR{each node $u \in \mathcal{V}$}
        \STATE $h_u^{agg} \leftarrow \frac{1}{|\mathcal{N}(u)|} \sum_{v \in \mathcal{N}(u)} v_v^{(L)}$
        \STATE $h_u^{new} \leftarrow \frac{1}{2}(v_u^{(L)} + h_u^{agg})$
        \STATE $\hat{y}_u \leftarrow \sigma(\text{MLP}(h_u^{new}))$
    \ENDFOR
    \STATE Compute $\mathcal{L}_{BCE}$ with $\text{pos\_weight} = 1.5 / \text{fraud ratio}$
    \STATE Update $\Theta$ via AdamW with cosine LR schedule
\ENDFOR
\end{algorithmic}
\end{algorithm}

\vspace{1em}

\begin{algorithm}[H]
\caption{SplitGNN Training Procedure}
\begin{algorithmic}[1]
\REQUIRE Graph $\mathcal{G}=(\mathcal{V},\mathcal{E},X)$, labels $\{y_u\}$, wavelet order $C$, split point $K$, edge weight $\gamma_e$
\ENSURE Trained model parameters $\Theta^*$
\STATE Initialize node embeddings $\{h_u\}$ from features $X$
\FOR{epoch $= 1, 2, \dots$ (with early stopping)}
    \FOR{each edge $(u,v) \in \mathcal{E}$}
        \STATE $\psi_{uv} \leftarrow \tanh(w^\top [h_u \parallel h_v \parallel (h_u - h_v)])$
    \ENDFOR
    \STATE $\mathcal{E}^+ \leftarrow \{(u,v) : \psi_{uv} > 0\}$; \quad $\mathcal{E}^- \leftarrow \{(u,v) : \psi_{uv} < 0\}$
    \STATE Construct $A^+, A^-$; compute Laplacians $L^+, L^-$ via CachedLaplacian
    \FOR{each node $u \in \mathcal{V}$}
        \STATE $\hat{h}_u^{(+)} \leftarrow \sum_{j=0}^{C-K} W_j^{(+)} g_j(L^+) h_u$
        \STATE $\hat{h}_u^{(-)} \leftarrow \sum_{j=\max(0,K)}^C W_j^{(-)} g_j(L^-) h_u$
        \STATE $\hat{h}_u \leftarrow \hat{h}_u^{(+)} + \hat{h}_u^{(-)}$
        \STATE Update $h_u \leftarrow \hat{h}_u$ \hfill (residual connection)
    \ENDFOR
    \STATE $\mathcal{L}_{node} \leftarrow \text{FocalLoss}(\hat{y}_u, y_u; \alpha=0.25, \gamma=2.0)$
    \STATE $\mathcal{L}_{edge} \leftarrow \frac{1}{|\mathcal{E}|} \sum_{(u,v) \in \mathcal{E}} \max(0, 1 - y_{uv}\psi_{uv})$
    \STATE $\mathcal{L} \leftarrow \mathcal{L}_{node} + \gamma_e \cdot \mathcal{L}_{edge}$
    \STATE Update $\Theta$ via full-batch gradient descent
\ENDFOR
\FOR{$t = 1$ to $5$}
    \STATE Retrain with $lr \leftarrow lr \times 0.1$
\ENDFOR
\end{algorithmic}
\end{algorithm}

\newpage

\section*{Appendix D: Extended Hyperparameter Configurations}

\subsection*{D.1 Split GNN: Top 5 Configurations}

The grid search evaluated interactions between the edge loss weight ($\gamma_e$), wavelet order ($C$), and split point ($K$). The top 5 stable configurations are listed below:

\begin{table}[H]
\centering
\caption{Top 5 Stable Configurations for Split GNN}
\vspace{0.5em}
\begin{tabular}{ccccc}
\toprule
\textbf{$\gamma_e$} & \textbf{$C$} & \textbf{$K$} & \textbf{Test ROC AUC} & \textbf{TPR @ 5\% FPR} \\ 
\midrule
1.0 & 1 &  0 & 0.8806 & 50.6\% \\
0.4 & 3 & -1 & 0.8804 & 50.1\% \\
0.8 & 1 &  0 & 0.8804 & 50.2\% \\
0.6 & 3 &  1 & 0.8803 & 49.6\% \\
0.6 & 1 &  0 & 0.8801 & 50.8\% \\
\bottomrule
\end{tabular}
\end{table}

\vspace{1.5em}

\subsection*{D.2 Quantum GNN (QGNN): Circuit Configuration Sweep}

Performance evaluation across quantum feature space dimensions (qubits) and Variational Quantum Circuit (VQC) depths.

\begin{table}[H]
\centering
\caption{QGNN Performance by Circuit Depth and Qubit Configuration}
\vspace{0.5em}
\begin{tabular}{lccccc}
\toprule
\textbf{Configuration} & \textbf{ROC AUC} & \textbf{TPR @ 5\% FPR} & \textbf{True Positives} & \textbf{False Negatives} & \textbf{Approx. Time} \\ 
\midrule
16Q, 2L & 0.8769 & 49.24\% & 1,417 & 1,461 & $\sim$75m \\
16Q, 1L & 0.8731 & 49.13\% & 1,414 & 1,464 & $\sim$70m \\
6Q, 1L  & 0.8602 & 46.04\% & 1,325 & 1,553 & $\sim$37m \\
6Q, 2L  & 0.8574 & 44.82\% & 1,290 & 1,588 & $\sim$38m \\
\bottomrule
\end{tabular}
\end{table}

\end{document}

```